% =====================================================================
%  Analiza istniejących rozwiązań i technologii — przegląd tego, co już
%  istnieje (biblioteki, narzędzia, podobne prace) z uzasadnieniem
%  wyboru technologii.
%
%  Demonstration deck: listing z kolorowaniem składni (minted, tu Python).
%  minted wymaga kompilacji z -shell-escape i zainstalowanego Pygments —
%  compile.sh załatwia jedno, a pip install Pygments drugie.
% =====================================================================
\section{Analiza istniejących rozwiązań i technologii}
Dzięki swojej prostocie, biblioteka stanowi standard w badaniach dotyczących
głębokiego uczenia oraz w zastosowaniach komercyjnych \cite{Paszke2019pytorch, Paszke2017, Lawson2024}.

\begin{minted}[
    frame=lines,        % Ramka na górze i na dole
    framesep=2mm,       % Odstęp ramki od kodu
    baselinestretch=1.2,% Interlinia dla czytelności
    fontsize=\small,    % Nieco mniejsza czcionka
    linenos             % Numery linii
]
{python}
import torch

# Operacje wykonywane natychmiast, graf budowany dynamicznie
x = torch.tensor([1.0, 2.0, 3.0], requires_grad=True)
y = x.pow(2)        # y.grad_fn = <PowBackward0>
z = y.sum()         # z.grad_fn = <SumBackward0>

print(z.item())     # Wartość dostępna od razu: 14.0

# Propagacja wsteczna przez zbudowany graf
z.backward()
print(x.grad)       # tensor([2., 4., 6.])
\end{minted}
